% =============================================================================
% Cover
% =============================================================================
\hypersetup{pageanchor=false}
\begin{titlepage}
  \thispagestyle{empty}
  \begin{tikzpicture}[remember picture,overlay]
    \fill[PrimaryDark] (current page.north west) rectangle ([yshift=-11.2cm]current page.north east);
    \fill[Primary] ([yshift=-11.2cm]current page.north west) rectangle ([yshift=-11.38cm]current page.north east);
    \begin{scope}[shift={([xshift=-1.6cm,yshift=-2.0cm]current page.north east)}]
      \foreach \i/\w in {0/3.4,1/4.3,2/5.2,3/6.1,4/7.0,5/7.9} {
        \fill[white,opacity=0.045+0.018*\i,rounded corners=1pt]
          (-\w,-1.05*\i-0.72) rectangle (0,-1.05*\i);
      }
    \end{scope}
    \node[anchor=north west,text=white,opacity=.78,font=\sffamily\small]
      at ([xshift=2.55cm,yshift=-2.15cm]current page.north west)
      {TECHNICAL WHITE PAPER \;·\; REFERENCE ARCHITECTURE};
    \draw[white,opacity=.55,line width=.6pt]
      ([xshift=2.55cm,yshift=-2.75cm]current page.north west) -- ++(1.8cm,0);
    \node[anchor=north west,text=white,font=\sffamily\bfseries\fontsize{34}{42}\selectfont,
          text width=15.5cm,align=left,inner sep=0pt]
      at ([xshift=2.55cm,yshift=-3.35cm]current page.north west)
      {Enterprise Agent\\Hosting Platform};
    \node[anchor=north west,text=white,opacity=.88,font=\sffamily\fontsize{15.5}{21}\selectfont,inner sep=0pt]
      at ([xshift=2.55cm,yshift=-7.45cm]current page.north west)
      {A Reference Architecture};
    \node[anchor=north west,text=white,opacity=.70,font=\sffamily\fontsize{10.5}{15}\selectfont,
          text width=13.6cm,align=left,inner sep=0pt]
      at ([xshift=2.55cm,yshift=-8.6cm]current page.north west)
      {Own the control plane. Rent the compute. Decouple through protocols.\\
       An architecture for hosting AI agents inside an enterprise without handing control to any single vendor.};
  \end{tikzpicture}

  \vspace*{10.3cm}

  \noindent
  \begin{minipage}{\textwidth}
    {\sffamily\bfseries\small\color{Primary}THE CORE PROPOSITION\par}
    \vspace{0.35em}
    {\color{Ink}\fontsize{12.5}{18}\selectfont
    A cloud vendor can host the machines. It cannot answer, on the enterprise's behalf, the questions that matter:
    \emph{who may run, what may they do, what did they do, what did it cost, and what happened when something went wrong}.
    The job of an enterprise agent platform is to put uncertain, model-driven behavior inside a boundary that is
    deterministic, verifiable, and portable.\par}
  \end{minipage}

  \vspace{1.4cm}

  \noindent
  \begin{tabularx}{\textwidth}{@{}YYY@{}}
    {\sffamily\bfseries\color{PrimaryDark}Own the control plane} &
    {\sffamily\bfseries\color{PrimaryDark}Rent the compute} &
    {\sffamily\bfseries\color{PrimaryDark}Decouple through protocols}\\[0.2em]
    {\small\color{BodyText}Identity, authorization, the ledger, task state and the interface contracts stay with the enterprise. They \emph{are} control.} &
    {\small\color{BodyText}Runtimes, sandboxes, object storage and dashboards are rented by usage. Any vendor that meets the contract will do.} &
    {\small\color{BodyText}The platform binds to a protocol, not to a framework. Any agent framework joins through an adapter; the contract never moves.}
  \end{tabularx}

  \vfill

  \noindent{\color{Rule}\rule{\textwidth}{0.6pt}}\par
  \vspace{0.55cm}
  \noindent
  \begin{tabularx}{\textwidth}{@{}L{2.4cm}L{6.2cm}L{2.4cm}Y@{}}
    {\sffamily\small\color{Muted}Author} & {\sffamily\bfseries\color{Ink}Yuting} &
    {\sffamily\small\color{Muted}Version} & {\sffamily\bfseries\color{Ink}v1.0}\\[0.35em]
    {\sffamily\small\color{Muted}Date} & {\sffamily\color{Ink}September 2026} &
    {\sffamily\small\color{Muted}Status} & {\sffamily\color{Ink}Reference architecture and design proposal}
  \end{tabularx}
  \vspace{0.45cm}

  \noindent{\small\color{Muted}
    This document describes a target architecture, its design constraints and a phased delivery boundary. It does not
    claim an existing production capability, measured performance, or any commercial service level.\par}
\end{titlepage}

\hypersetup{pageanchor=true}
\pagenumbering{Roman}

% =============================================================================
% Abstract
% =============================================================================
\section*{Abstract}

AI agents inside enterprises are moving from scattered experiments to managed hosting. What decides whether an agent can enter production is rarely the model itself. It is the set of questions around the model: whose identity does it act under, which systems may it touch, where do its credentials live, how far did a task get, what did it cost, how is an incident reconstructed, and what does it take to leave a vendor. When those capabilities are built on top of a single cloud runtime, the enterprise gains speed and loses control over its identities, policies, secrets and history.

This paper proposes a general reference architecture for an enterprise agent hosting platform. It is a design proposal, not an implementation report. Four decisions shape it. The enterprise keeps its own control plane---identity, authorization, the ledger, task state and the interface contracts---while renting machines, sandboxes, object storage and dashboards by usage. The platform binds to a protocol rather than to any agent framework. Every action an agent takes against the outside world---calling a model, calling a tool, executing code---passes through a governed exit that the platform owns. And every fact about a run is written to a ledger the enterprise holds.

On top of the managed runtimes, workload identity, credential vaults, gateways, sandboxes and durable execution that the industry already provides, the paper adds a set of enterprise-side mechanisms: a four-axis capability manifest that describes an agent along independent dimensions; a minimal container contract in which the platform, not the container, assigns event sequence numbers; a permission-intersection rule that covers user delegation, scheduled principals and structural denials; a three-stage kill switch that still works when the cloud vendor is unreachable; ``effectively-once'' side-effect semantics built from checkpoints, an idempotency ledger and archived results; and a hash-chained enterprise ledger that supports compliant deletion.

{\small\sffamily\color{Muted}
  \textbf{\color{Ink}Keywords:} enterprise agents; agent hosting platform; control plane; workload identity;
  tool governance; unified ledger; sandbox; cloud portability\par}

\subsection*{How to read this document}

The document is long, and different readers need different parts of it. Readers new to agent platforms should start with the glossary in Appendix~\ref{app:glossary}, then read Sections~\ref{sec:exec}, \ref{sec:problem} and \ref{sec:agent-def}, and then jump to Section~\ref{sec:lifecycle}, which walks a single request through every mechanism in the paper.

\begin{tabularx}{\textwidth}{@{}L{3.6cm}Y@{}}
  \toprule
  \rowcolor{PrimaryLight}
  \textbf{Reader} & \textbf{Suggested route}\\
  \midrule
  Executives and decision makers & Part I (the case and the object), Section~\ref{sec:overview} on the overall architecture, Section~\ref{sec:mvp} on delivery, Section~\ref{sec:conclusion}\\
  Platform architects & The whole document, with emphasis on the ten invariants, the three interface layers, runtime semantics, identity and authorization, and the architecture decisions\\
  Implementation teams & Parts II and III plus Appendices~\ref{app:l2} and \ref{app:checklist}: interfaces, runtime, identity, ledger and acceptance\\
  Security and audit reviewers & Section~\ref{sec:threat} (threat model), Sections~\ref{sec:iam}--\ref{sec:ledger}, Section~\ref{sec:risks} (residual risks)\\
  First-time readers & Appendix~\ref{app:glossary} $\rightarrow$ Sections~\ref{sec:exec}--\ref{sec:agent-def} $\rightarrow$ Section~\ref{sec:lifecycle}\\
  \bottomrule
\end{tabularx}

\newpage
\begingroup
\footnotesize
\setlength{\parskip}{0pt}
\renewcommand{\baselinestretch}{0.92}\selectfont
\tableofcontents
\vspace{1.2em}
\listoffigures
\vspace{1.2em}
\listoftables
\endgroup
\newpage
\pagenumbering{arabic}
